% GENERATED FILE. Edit canonical lessons, then run npm run generate:books.
% canonical-source-hash: fnv1a64:cf9cefbbc6dd95de
% canonical-lessons: MW-C134-makkhan, MW-C134-chhachh, MW-C134-lassi, MW-C134-kachori, MW-C134-bhujiya

\chapter{Butter, Buttermilk, Lassi, Kachori, Bhujia}
\label{ch:mw-butter-buttermilk-lassi-kachori-bhujia}
\hlchaptermodality{134}

\begin{chapteropening}
\textbf{By the end of this chapter:} \emph{I can say butter, buttermilk, lassi, a kachori and bhujia.}

Complete the last lesson of chapter 134: I can say butter, buttermilk, lassi, a kachori and bhujia.
\end{chapteropening}

\section[\texorpdfstring{\mw{मक्खण} (makkhaṇ)}{मक्खण (makkhaṇ)}]{\mw{मक्खण} (makkhaṇ) — butter}

\label{lesson:MW-C134-makkhan}



\begin{quote}
Before the new one: say the Marwadi for sesame, then the Marwadi for jaggery.
\end{quote}

\subsection*{You'll want to know: \mw{मक्खण}}
\textbf{\mw{मक्खण}} — \emph{makkhaṇ} — \textquotedblleft{}butter\textquotedblright{}.

\begin{cousinweb}
\textbf{\mw{मक्खण}} (\emph{makkhaṇ}) is a tadbhava, a word inherited from Sanskrit and reshaped by sound change. It comes from \emph{mrakṣaṇa}, \textquotedblleft{}rubbing in; ointment, oil\textquotedblright{}. CDIAL gives no Marwari form in this entry (CDIAL 10378); the match is through Hindi \emph{mākhan} and Gujarati \emph{mākhaṇ}.
\end{cousinweb}

\begin{grammarlens}[title={What you've built}]
One more word for this chapter.
\end{grammarlens}

\subsection*{Your turn}
\begin{itemize}
\raggedright
  \item \emph{Say it:} \emph{makkhaṇ}
  \item \emph{Say it:} \emph{makkhaṇ}, once more
  \item \emph{Say it:} say \emph{makkhaṇ}
  \item \emph{Recall:} say the Marwadi for sesame, then the Marwadi for jaggery, then say \emph{makkhaṇ} again
\end{itemize}

\begin{tcolorbox}[breakable,colback=teal!4,colframe=teal!35!black,title={Before you move on}]
What does \mw{मक्खण} mean? (\textquotedblleft{}Butter\textquotedblright{}.) Say it once more.
\end{tcolorbox}
\section[\texorpdfstring{\mw{छाछ} (chhāchh)}{छाछ (chhāchh)}]{\mw{छाछ} (chhāchh) — buttermilk}

\label{lesson:MW-C134-chhachh}



\begin{quote}
Before the new one: say the Marwadi for jaggery, then the Marwadi for butter.
\end{quote}

\subsection*{You'll want to know: \mw{छाछ}}
\textbf{\mw{छाछ}} — \emph{chhāchh} — \textquotedblleft{}buttermilk\textquotedblright{}.

\begin{cousinweb}
\textbf{\mw{छाछ}} (\emph{chhāchh}) is a tadbhava, a word inherited and reshaped by sound change. It comes from \emph{*chācchī}, \textquotedblleft{}buttermilk\textquotedblright{}, a form CDIAL reconstructs from its descendants. CDIAL lists Marwari \emph{chāchī} under it (CDIAL 5012).
\end{cousinweb}

\begin{grammarlens}[title={What you've built}]
One more word for this chapter.
\end{grammarlens}

\subsection*{Your turn}
\begin{itemize}
\raggedright
  \item \emph{Say it:} \emph{chhāchh}
  \item \emph{Say it:} \emph{chhāchh}, once more
  \item \emph{Say it:} say \emph{chhāchh}
  \item \emph{Recall:} say the Marwadi for jaggery, then the Marwadi for butter, then say \emph{chhāchh} again
\end{itemize}

\begin{tcolorbox}[breakable,colback=teal!4,colframe=teal!35!black,title={Before you move on}]
What does \mw{छाछ} mean? (\textquotedblleft{}Buttermilk\textquotedblright{}.) Say it once more.
\end{tcolorbox}
\section[\texorpdfstring{\mw{लस्सी} (lassī)}{लस्सी (lassī)}]{\mw{लस्सी} (lassī) — lassi}

\label{lesson:MW-C134-lassi}



\begin{quote}
Before the new one: say the Marwadi for butter, then the Marwadi for buttermilk.
\end{quote}

\subsection*{You'll want to know: \mw{लस्सी}}
\textbf{\mw{लस्सी}} — \emph{lassī} — \textquotedblleft{}lassi\textquotedblright{}.

\begin{cousinweb}
CDIAL reconstructs \emph{*lassī}, \textquotedblleft{}buttermilk\textquotedblright{}, from Punjabi and Sindhi forms (CDIAL 10998). It lists no Marwari or Hindi form, and has no older source for it.
\end{cousinweb}

\begin{grammarlens}[title={What you've built}]
One more word for this chapter.
\end{grammarlens}

\subsection*{Your turn}
\begin{itemize}
\raggedright
  \item \emph{Say it:} \emph{lassī}
  \item \emph{Say it:} \emph{lassī}, once more
  \item \emph{Say it:} say \emph{lassī}
  \item \emph{Recall:} say the Marwadi for butter, then the Marwadi for buttermilk, then say \emph{lassī} again
\end{itemize}

\begin{tcolorbox}[breakable,colback=teal!4,colframe=teal!35!black,title={Before you move on}]
What does \mw{लस्सी} mean? (\textquotedblleft{}Lassi\textquotedblright{}.) Say it once more.
\end{tcolorbox}
\section[\texorpdfstring{\mw{कचोरी} (kachorī)}{कचोरी (kachorī)}]{\mw{कचोरी} (kachorī) — a kachori}

\label{lesson:MW-C134-kachori}



\begin{quote}
Before the new one: say the Marwadi for buttermilk, then the Marwadi for lassi.
\end{quote}

\subsection*{You'll want to know: \mw{कचोरी}}
\textbf{\mw{कचोरी}} — \emph{kachorī} — \textquotedblleft{}a kachori\textquotedblright{}.

\begin{cousinweb}
\textbf{\mw{कचोरी}} (\emph{kachorī}) is a tadbhava, a word inherited and reshaped by sound change. It comes from \emph{*kaccapūra}, \textquotedblleft{}a wheaten cake\textquotedblright{}, a form CDIAL reconstructs from its descendants. CDIAL gives no Marwari form in this entry (CDIAL 2614); the match is through Hindi \emph{kacaurī} and Gujarati \emph{kacorī}. CDIAL builds it on \emph{pūra}, \textquotedblleft{}cake\textquotedblright{}, with a first part it queries.
\end{cousinweb}

\begin{grammarlens}[title={What you've built}]
One more word for this chapter.
\end{grammarlens}

\subsection*{Your turn}
\begin{itemize}
\raggedright
  \item \emph{Say it:} \emph{kachorī}
  \item \emph{Say it:} \emph{kachorī}, once more
  \item \emph{Say it:} say \emph{kachorī}
  \item \emph{Recall:} say the Marwadi for buttermilk, then the Marwadi for lassi, then say \emph{kachorī} again
\end{itemize}

\begin{tcolorbox}[breakable,colback=teal!4,colframe=teal!35!black,title={Before you move on}]
What does \mw{कचोरी} mean? (\textquotedblleft{}A kachori\textquotedblright{}.) Say it once more.
\end{tcolorbox}
\section[\texorpdfstring{\mw{भुजिया} (bhujiyā)}{भुजिया (bhujiyā)}]{\mw{भुजिया} (bhujiyā) — bhujia}

\label{lesson:MW-C134-bhujiya}



\begin{quote}
Before the new one: say the Marwadi for lassi, then the Marwadi for a kachori.
\end{quote}

\subsection*{You'll want to know: \mw{भुजिया}}
\textbf{\mw{भुजिया}} — \emph{bhujiyā} — \textquotedblleft{}bhujia\textquotedblright{}.

\begin{cousinweb}
\textbf{\mw{भुजिया}} (\emph{bhujiyā}) is a tadbhava, a word inherited and reshaped by sound change. It comes from \emph{*bhṛjjita}, \textquotedblleft{}fried, parched\textquotedblright{}, a form CDIAL reconstructs from its descendants. CDIAL gives no Marwari form in this entry (CDIAL 9584); the match is through Hindi \emph{bhujiyā}. It is built on the root \emph{bhrajj}, \textquotedblleft{}to fry\textquotedblright{}.
\end{cousinweb}

\begin{grammarlens}[title={What you've built}]
One more word for this chapter.
\end{grammarlens}

\subsection*{Your turn}
\begin{itemize}
\raggedright
  \item \emph{Say it:} \emph{bhujiyā}
  \item \emph{Say it:} \emph{bhujiyā}, once more
  \item \emph{Say it:} say \emph{bhujiyā}
  \item \emph{Recall:} say the Marwadi for lassi, then the Marwadi for a kachori, then say \emph{bhujiyā} again
\end{itemize}

\begin{tcolorbox}[breakable,colback=teal!4,colframe=teal!35!black,title={Before you move on}]
What does \mw{भुजिया} mean? (\textquotedblleft{}Bhujia\textquotedblright{}.) Say it once more.
\end{tcolorbox}
